% =====================================================================
%  Psicoestadística Inferencial — Unidad 2 (Probabilidad)
%  Preámbulo común a todas las salidas.
%
%  Un solo código fuente produce cuatro documentos distintos, según dos
%  interruptores que se fijan ANTES de % =====================================================================
%  Psicoestadística Inferencial — Unidad 2 (Probabilidad)
%  Preámbulo común a todas las salidas.
%
%  Un solo código fuente produce cuatro documentos distintos, según dos
%  interruptores que se fijan ANTES de % =====================================================================
%  Psicoestadística Inferencial — Unidad 2 (Probabilidad)
%  Preámbulo común a todas las salidas.
%
%  Un solo código fuente produce cuatro documentos distintos, según dos
%  interruptores que se fijan ANTES de % =====================================================================
%  Psicoestadística Inferencial — Unidad 2 (Probabilidad)
%  Preámbulo común a todas las salidas.
%
%  Un solo código fuente produce cuatro documentos distintos, según dos
%  interruptores que se fijan ANTES de \input{preambulo}:
%
%     \docentetrue    incluye la solución paso a paso
%     \docentefalse   sólo enunciados (los estudiantes)
%     \espaciotrue    deja espacio en blanco para resolver a mano
%     \espaciofalse   compacto, para leer en pantalla
%
%  Así los enunciados nunca se desincronizan de sus soluciones.
% =====================================================================

% `lmodern` da tipografías vectoriales: sin esto microtype aborta la
% compilación («auto expansion is only possible with scalable fonts»).
\usepackage{lmodern}
\usepackage[T1]{fontenc}
\usepackage[utf8]{inputenc}
\usepackage[spanish,es-noquoting,es-nodecimaldot]{babel}
\usepackage{amsmath,amssymb}
\usepackage[table]{xcolor}
\usepackage{booktabs}
\usepackage{enumitem}
\usepackage{fancyhdr}
\usepackage[most]{tcolorbox}
\usepackage{lastpage}
\usepackage{microtype}

% --- Colores: los mismos bloques temáticos que la Aula interactiva ---
\definecolor{tinta}{HTML}{0F172A}
\definecolor{apagado}{HTML}{64748B}
\definecolor{borde}{HTML}{CBD5E1}
\definecolor{suave}{HTML}{F1F5F9}
\definecolor{nivelcero}{HTML}{64748B}   % reconocer
\definecolor{niveluno}{HTML}{2563EB}    % contar
\definecolor{niveldos}{HTML}{4F46E5}    % una fórmula
\definecolor{niveltres}{HTML}{D97706}   % datos reales
\definecolor{nivelcuatro}{HTML}{059669} % decidir

\color{tinta}

% --- Encabezado y pie ---
\pagestyle{fancy}
\fancyhf{}
\renewcommand{\headrulewidth}{0.4pt}
\renewcommand{\footrulewidth}{0pt}
\fancyhead[L]{\footnotesize\textsc{Psicoestadística Inferencial}}
\fancyhead[R]{\footnotesize\textsc{Unidad 2 · Probabilidad}}
\fancyfoot[C]{\footnotesize\color{apagado}\thepage\ de \pageref{LastPage}}

% --- Contadores ---
\newcounter{numejercicio}
\setcounter{numejercicio}{0}

% --- Nombres de los cinco niveles ---
\newcommand{\nombrenivel}[1]{%
  \ifcase#1 Reconocer\or Contar\or Una fórmula\or Datos reales\or Decidir\fi}
\newcommand{\colornivel}[1]{%
  \ifcase#1 nivelcero\or niveluno\or niveldos\or niveltres\or nivelcuatro\fi}

% =====================================================================
%  \ejercicio{nivel}{enunciado}
% =====================================================================
\newcommand{\ejercicio}[2]{%
  \stepcounter{numejercicio}%
  \par\vspace{1.1em}%
  \noindent
  {\color{\colornivel{#1}}\textbf{\thenumejercicio.}}\quad
  \raisebox{0.12em}{\footnotesize\colorbox{\colornivel{#1}}{\textcolor{white}{\bfseries\,N#1 · \nombrenivel{#1}\,}}}%
  \par\vspace{0.35em}%
  \noindent #2\par
}

% --- Espacio para resolver a mano (sólo si \espaciotrue) ---
\newcommand{\espacioresolver}[1][3.2cm]{%
  \ifespacio\vspace{#1}\fi}

% =====================================================================
%  \solucion{...}   paso a paso — sólo en la versión docente
%  \respuesta{...}  resultado corto — va al final en la del estudiante
% =====================================================================
\newcommand{\solucion}[1]{%
  \ifdocente
    \par\vspace{0.4em}
    \begin{tcolorbox}[
      colback=suave, colframe=borde, boxrule=0.5pt, arc=3pt,
      left=8pt, right=8pt, top=6pt, bottom=6pt, breakable]
      {\footnotesize\textsc{\color{apagado}Solución}}\par\vspace{0.3em}#1
    \end{tcolorbox}
  \fi}

% El resultado corto se acumula en una macro a medida que aparecen los
% ejercicios, y se imprime al final como clave de respuestas. Se acumula en
% memoria en vez de escribirse a un archivo auxiliar porque escribir expande
% el contenido, y cualquier fórmula matemática de una respuesta se rompería.
\usepackage{etoolbox}
\newcommand{\claves}{}
\newcommand{\respuesta}[1]{%
  \edef\numactual{\thenumejercicio}%
  \expandafter\gappto\expandafter\claves\expandafter{%
    \expandafter\item\expandafter[\numactual.]}%
  \gappto\claves{#1}%
}

\newcommand{\imprimirclave}{%
  \ifdocente\else
    \clearpage
    \section*{Respuestas}
    \noindent{\color{apagado}\small Sólo el resultado. El desarrollo se
    trabaja en clase — si tu número no coincide, el error está en algún paso
    intermedio y vale la pena encontrarlo antes de mirar la solución.}
    \par\vspace{0.8em}
    \begin{multicols}{2}\raggedright
    \begin{description}[leftmargin=2.8em,style=sameline,itemsep=0.4em]
      \claves
    \end{description}
    \end{multicols}
  \fi}

\usepackage{multicol}

% --- Encabezado de apartado ---
\newcommand{\apartado}[2]{%
  \clearpage
  \noindent{\color{apagado}\footnotesize\textsc{Apartado}}\par
  \noindent{\LARGE\bfseries #1\quad #2}\par
  \vspace{0.2em}\noindent\rule{\linewidth}{0.8pt}\par\vspace{0.3em}}

% --- Nota al docente / recordatorio de fórmula ---
\newcommand{\recordar}[1]{%
  \par\vspace{0.3em}
  \noindent\begin{tcolorbox}[
    colback=white, colframe=borde, boxrule=0.5pt, arc=3pt,
    left=8pt, right=8pt, top=5pt, bottom=5pt]
    {\footnotesize #1}
  \end{tcolorbox}}
:
%
%     \docentetrue    incluye la solución paso a paso
%     \docentefalse   sólo enunciados (los estudiantes)
%     \espaciotrue    deja espacio en blanco para resolver a mano
%     \espaciofalse   compacto, para leer en pantalla
%
%  Así los enunciados nunca se desincronizan de sus soluciones.
% =====================================================================

% `lmodern` da tipografías vectoriales: sin esto microtype aborta la
% compilación («auto expansion is only possible with scalable fonts»).
\usepackage{lmodern}
\usepackage[T1]{fontenc}
\usepackage[utf8]{inputenc}
\usepackage[spanish,es-noquoting,es-nodecimaldot]{babel}
\usepackage{amsmath,amssymb}
\usepackage[table]{xcolor}
\usepackage{booktabs}
\usepackage{enumitem}
\usepackage{fancyhdr}
\usepackage[most]{tcolorbox}
\usepackage{lastpage}
\usepackage{microtype}

% --- Colores: los mismos bloques temáticos que la Aula interactiva ---
\definecolor{tinta}{HTML}{0F172A}
\definecolor{apagado}{HTML}{64748B}
\definecolor{borde}{HTML}{CBD5E1}
\definecolor{suave}{HTML}{F1F5F9}
\definecolor{nivelcero}{HTML}{64748B}   % reconocer
\definecolor{niveluno}{HTML}{2563EB}    % contar
\definecolor{niveldos}{HTML}{4F46E5}    % una fórmula
\definecolor{niveltres}{HTML}{D97706}   % datos reales
\definecolor{nivelcuatro}{HTML}{059669} % decidir

\color{tinta}

% --- Encabezado y pie ---
\pagestyle{fancy}
\fancyhf{}
\renewcommand{\headrulewidth}{0.4pt}
\renewcommand{\footrulewidth}{0pt}
\fancyhead[L]{\footnotesize\textsc{Psicoestadística Inferencial}}
\fancyhead[R]{\footnotesize\textsc{Unidad 2 · Probabilidad}}
\fancyfoot[C]{\footnotesize\color{apagado}\thepage\ de \pageref{LastPage}}

% --- Contadores ---
\newcounter{numejercicio}
\setcounter{numejercicio}{0}

% --- Nombres de los cinco niveles ---
\newcommand{\nombrenivel}[1]{%
  \ifcase#1 Reconocer\or Contar\or Una fórmula\or Datos reales\or Decidir\fi}
\newcommand{\colornivel}[1]{%
  \ifcase#1 nivelcero\or niveluno\or niveldos\or niveltres\or nivelcuatro\fi}

% =====================================================================
%  \ejercicio{nivel}{enunciado}
% =====================================================================
\newcommand{\ejercicio}[2]{%
  \stepcounter{numejercicio}%
  \par\vspace{1.1em}%
  \noindent
  {\color{\colornivel{#1}}\textbf{\thenumejercicio.}}\quad
  \raisebox{0.12em}{\footnotesize\colorbox{\colornivel{#1}}{\textcolor{white}{\bfseries\,N#1 · \nombrenivel{#1}\,}}}%
  \par\vspace{0.35em}%
  \noindent #2\par
}

% --- Espacio para resolver a mano (sólo si \espaciotrue) ---
\newcommand{\espacioresolver}[1][3.2cm]{%
  \ifespacio\vspace{#1}\fi}

% =====================================================================
%  \solucion{...}   paso a paso — sólo en la versión docente
%  \respuesta{...}  resultado corto — va al final en la del estudiante
% =====================================================================
\newcommand{\solucion}[1]{%
  \ifdocente
    \par\vspace{0.4em}
    \begin{tcolorbox}[
      colback=suave, colframe=borde, boxrule=0.5pt, arc=3pt,
      left=8pt, right=8pt, top=6pt, bottom=6pt, breakable]
      {\footnotesize\textsc{\color{apagado}Solución}}\par\vspace{0.3em}#1
    \end{tcolorbox}
  \fi}

% El resultado corto se acumula en una macro a medida que aparecen los
% ejercicios, y se imprime al final como clave de respuestas. Se acumula en
% memoria en vez de escribirse a un archivo auxiliar porque escribir expande
% el contenido, y cualquier fórmula matemática de una respuesta se rompería.
\usepackage{etoolbox}
\newcommand{\claves}{}
\newcommand{\respuesta}[1]{%
  \edef\numactual{\thenumejercicio}%
  \expandafter\gappto\expandafter\claves\expandafter{%
    \expandafter\item\expandafter[\numactual.]}%
  \gappto\claves{#1}%
}

\newcommand{\imprimirclave}{%
  \ifdocente\else
    \clearpage
    \section*{Respuestas}
    \noindent{\color{apagado}\small Sólo el resultado. El desarrollo se
    trabaja en clase — si tu número no coincide, el error está en algún paso
    intermedio y vale la pena encontrarlo antes de mirar la solución.}
    \par\vspace{0.8em}
    \begin{multicols}{2}\raggedright
    \begin{description}[leftmargin=2.8em,style=sameline,itemsep=0.4em]
      \claves
    \end{description}
    \end{multicols}
  \fi}

\usepackage{multicol}

% --- Encabezado de apartado ---
\newcommand{\apartado}[2]{%
  \clearpage
  \noindent{\color{apagado}\footnotesize\textsc{Apartado}}\par
  \noindent{\LARGE\bfseries #1\quad #2}\par
  \vspace{0.2em}\noindent\rule{\linewidth}{0.8pt}\par\vspace{0.3em}}

% --- Nota al docente / recordatorio de fórmula ---
\newcommand{\recordar}[1]{%
  \par\vspace{0.3em}
  \noindent\begin{tcolorbox}[
    colback=white, colframe=borde, boxrule=0.5pt, arc=3pt,
    left=8pt, right=8pt, top=5pt, bottom=5pt]
    {\footnotesize #1}
  \end{tcolorbox}}
:
%
%     \docentetrue    incluye la solución paso a paso
%     \docentefalse   sólo enunciados (los estudiantes)
%     \espaciotrue    deja espacio en blanco para resolver a mano
%     \espaciofalse   compacto, para leer en pantalla
%
%  Así los enunciados nunca se desincronizan de sus soluciones.
% =====================================================================

% `lmodern` da tipografías vectoriales: sin esto microtype aborta la
% compilación («auto expansion is only possible with scalable fonts»).
\usepackage{lmodern}
\usepackage[T1]{fontenc}
\usepackage[utf8]{inputenc}
\usepackage[spanish,es-noquoting,es-nodecimaldot]{babel}
\usepackage{amsmath,amssymb}
\usepackage[table]{xcolor}
\usepackage{booktabs}
\usepackage{enumitem}
\usepackage{fancyhdr}
\usepackage[most]{tcolorbox}
\usepackage{lastpage}
\usepackage{microtype}

% --- Colores: los mismos bloques temáticos que la Aula interactiva ---
\definecolor{tinta}{HTML}{0F172A}
\definecolor{apagado}{HTML}{64748B}
\definecolor{borde}{HTML}{CBD5E1}
\definecolor{suave}{HTML}{F1F5F9}
\definecolor{nivelcero}{HTML}{64748B}   % reconocer
\definecolor{niveluno}{HTML}{2563EB}    % contar
\definecolor{niveldos}{HTML}{4F46E5}    % una fórmula
\definecolor{niveltres}{HTML}{D97706}   % datos reales
\definecolor{nivelcuatro}{HTML}{059669} % decidir

\color{tinta}

% --- Encabezado y pie ---
\pagestyle{fancy}
\fancyhf{}
\renewcommand{\headrulewidth}{0.4pt}
\renewcommand{\footrulewidth}{0pt}
\fancyhead[L]{\footnotesize\textsc{Psicoestadística Inferencial}}
\fancyhead[R]{\footnotesize\textsc{Unidad 2 · Probabilidad}}
\fancyfoot[C]{\footnotesize\color{apagado}\thepage\ de \pageref{LastPage}}

% --- Contadores ---
\newcounter{numejercicio}
\setcounter{numejercicio}{0}

% --- Nombres de los cinco niveles ---
\newcommand{\nombrenivel}[1]{%
  \ifcase#1 Reconocer\or Contar\or Una fórmula\or Datos reales\or Decidir\fi}
\newcommand{\colornivel}[1]{%
  \ifcase#1 nivelcero\or niveluno\or niveldos\or niveltres\or nivelcuatro\fi}

% =====================================================================
%  \ejercicio{nivel}{enunciado}
% =====================================================================
\newcommand{\ejercicio}[2]{%
  \stepcounter{numejercicio}%
  \par\vspace{1.1em}%
  \noindent
  {\color{\colornivel{#1}}\textbf{\thenumejercicio.}}\quad
  \raisebox{0.12em}{\footnotesize\colorbox{\colornivel{#1}}{\textcolor{white}{\bfseries\,N#1 · \nombrenivel{#1}\,}}}%
  \par\vspace{0.35em}%
  \noindent #2\par
}

% --- Espacio para resolver a mano (sólo si \espaciotrue) ---
\newcommand{\espacioresolver}[1][3.2cm]{%
  \ifespacio\vspace{#1}\fi}

% =====================================================================
%  \solucion{...}   paso a paso — sólo en la versión docente
%  \respuesta{...}  resultado corto — va al final en la del estudiante
% =====================================================================
\newcommand{\solucion}[1]{%
  \ifdocente
    \par\vspace{0.4em}
    \begin{tcolorbox}[
      colback=suave, colframe=borde, boxrule=0.5pt, arc=3pt,
      left=8pt, right=8pt, top=6pt, bottom=6pt, breakable]
      {\footnotesize\textsc{\color{apagado}Solución}}\par\vspace{0.3em}#1
    \end{tcolorbox}
  \fi}

% El resultado corto se acumula en una macro a medida que aparecen los
% ejercicios, y se imprime al final como clave de respuestas. Se acumula en
% memoria en vez de escribirse a un archivo auxiliar porque escribir expande
% el contenido, y cualquier fórmula matemática de una respuesta se rompería.
\usepackage{etoolbox}
\newcommand{\claves}{}
\newcommand{\respuesta}[1]{%
  \edef\numactual{\thenumejercicio}%
  \expandafter\gappto\expandafter\claves\expandafter{%
    \expandafter\item\expandafter[\numactual.]}%
  \gappto\claves{#1}%
}

\newcommand{\imprimirclave}{%
  \ifdocente\else
    \clearpage
    \section*{Respuestas}
    \noindent{\color{apagado}\small Sólo el resultado. El desarrollo se
    trabaja en clase — si tu número no coincide, el error está en algún paso
    intermedio y vale la pena encontrarlo antes de mirar la solución.}
    \par\vspace{0.8em}
    \begin{multicols}{2}\raggedright
    \begin{description}[leftmargin=2.8em,style=sameline,itemsep=0.4em]
      \claves
    \end{description}
    \end{multicols}
  \fi}

\usepackage{multicol}

% --- Encabezado de apartado ---
\newcommand{\apartado}[2]{%
  \clearpage
  \noindent{\color{apagado}\footnotesize\textsc{Apartado}}\par
  \noindent{\LARGE\bfseries #1\quad #2}\par
  \vspace{0.2em}\noindent\rule{\linewidth}{0.8pt}\par\vspace{0.3em}}

% --- Nota al docente / recordatorio de fórmula ---
\newcommand{\recordar}[1]{%
  \par\vspace{0.3em}
  \noindent\begin{tcolorbox}[
    colback=white, colframe=borde, boxrule=0.5pt, arc=3pt,
    left=8pt, right=8pt, top=5pt, bottom=5pt]
    {\footnotesize #1}
  \end{tcolorbox}}
:
%
%     \docentetrue    incluye la solución paso a paso
%     \docentefalse   sólo enunciados (los estudiantes)
%     \espaciotrue    deja espacio en blanco para resolver a mano
%     \espaciofalse   compacto, para leer en pantalla
%
%  Así los enunciados nunca se desincronizan de sus soluciones.
% =====================================================================

% `lmodern` da tipografías vectoriales: sin esto microtype aborta la
% compilación («auto expansion is only possible with scalable fonts»).
\usepackage{lmodern}
\usepackage[T1]{fontenc}
\usepackage[utf8]{inputenc}
\usepackage[spanish,es-noquoting,es-nodecimaldot]{babel}
\usepackage{amsmath,amssymb}
\usepackage[table]{xcolor}
\usepackage{booktabs}
\usepackage{enumitem}
\usepackage{fancyhdr}
\usepackage[most]{tcolorbox}
\usepackage{lastpage}
\usepackage{microtype}

% --- Colores: los mismos bloques temáticos que la Aula interactiva ---
\definecolor{tinta}{HTML}{0F172A}
\definecolor{apagado}{HTML}{64748B}
\definecolor{borde}{HTML}{CBD5E1}
\definecolor{suave}{HTML}{F1F5F9}
\definecolor{nivelcero}{HTML}{64748B}   % reconocer
\definecolor{niveluno}{HTML}{2563EB}    % contar
\definecolor{niveldos}{HTML}{4F46E5}    % una fórmula
\definecolor{niveltres}{HTML}{D97706}   % datos reales
\definecolor{nivelcuatro}{HTML}{059669} % decidir

\color{tinta}

% --- Encabezado y pie ---
\pagestyle{fancy}
\fancyhf{}
\renewcommand{\headrulewidth}{0.4pt}
\renewcommand{\footrulewidth}{0pt}
\fancyhead[L]{\footnotesize\textsc{Psicoestadística Inferencial}}
\fancyhead[R]{\footnotesize\textsc{Unidad 2 · Probabilidad}}
\fancyfoot[C]{\footnotesize\color{apagado}\thepage\ de \pageref{LastPage}}

% --- Contadores ---
\newcounter{numejercicio}
\setcounter{numejercicio}{0}

% --- Nombres de los cinco niveles ---
\newcommand{\nombrenivel}[1]{%
  \ifcase#1 Reconocer\or Contar\or Una fórmula\or Datos reales\or Decidir\fi}
\newcommand{\colornivel}[1]{%
  \ifcase#1 nivelcero\or niveluno\or niveldos\or niveltres\or nivelcuatro\fi}

% =====================================================================
%  \ejercicio{nivel}{enunciado}
% =====================================================================
\newcommand{\ejercicio}[2]{%
  \stepcounter{numejercicio}%
  \par\vspace{1.1em}%
  \noindent
  {\color{\colornivel{#1}}\textbf{\thenumejercicio.}}\quad
  \raisebox{0.12em}{\footnotesize\colorbox{\colornivel{#1}}{\textcolor{white}{\bfseries\,N#1 · \nombrenivel{#1}\,}}}%
  \par\vspace{0.35em}%
  \noindent #2\par
}

% --- Espacio para resolver a mano (sólo si \espaciotrue) ---
\newcommand{\espacioresolver}[1][3.2cm]{%
  \ifespacio\vspace{#1}\fi}

% =====================================================================
%  \solucion{...}   paso a paso — sólo en la versión docente
%  \respuesta{...}  resultado corto — va al final en la del estudiante
% =====================================================================
\newcommand{\solucion}[1]{%
  \ifdocente
    \par\vspace{0.4em}
    \begin{tcolorbox}[
      colback=suave, colframe=borde, boxrule=0.5pt, arc=3pt,
      left=8pt, right=8pt, top=6pt, bottom=6pt, breakable]
      {\footnotesize\textsc{\color{apagado}Solución}}\par\vspace{0.3em}#1
    \end{tcolorbox}
  \fi}

% El resultado corto se acumula en una macro a medida que aparecen los
% ejercicios, y se imprime al final como clave de respuestas. Se acumula en
% memoria en vez de escribirse a un archivo auxiliar porque escribir expande
% el contenido, y cualquier fórmula matemática de una respuesta se rompería.
\usepackage{etoolbox}
\newcommand{\claves}{}
\newcommand{\respuesta}[1]{%
  \edef\numactual{\thenumejercicio}%
  \expandafter\gappto\expandafter\claves\expandafter{%
    \expandafter\item\expandafter[\numactual.]}%
  \gappto\claves{#1}%
}

\newcommand{\imprimirclave}{%
  \ifdocente\else
    \clearpage
    \section*{Respuestas}
    \noindent{\color{apagado}\small Sólo el resultado. El desarrollo se
    trabaja en clase — si tu número no coincide, el error está en algún paso
    intermedio y vale la pena encontrarlo antes de mirar la solución.}
    \par\vspace{0.8em}
    \begin{multicols}{2}\raggedright
    \begin{description}[leftmargin=2.8em,style=sameline,itemsep=0.4em]
      \claves
    \end{description}
    \end{multicols}
  \fi}

\usepackage{multicol}

% --- Encabezado de apartado ---
\newcommand{\apartado}[2]{%
  \clearpage
  \noindent{\color{apagado}\footnotesize\textsc{Apartado}}\par
  \noindent{\LARGE\bfseries #1\quad #2}\par
  \vspace{0.2em}\noindent\rule{\linewidth}{0.8pt}\par\vspace{0.3em}}

% --- Nota al docente / recordatorio de fórmula ---
\newcommand{\recordar}[1]{%
  \par\vspace{0.3em}
  \noindent\begin{tcolorbox}[
    colback=white, colframe=borde, boxrule=0.5pt, arc=3pt,
    left=8pt, right=8pt, top=5pt, bottom=5pt]
    {\footnotesize #1}
  \end{tcolorbox}}
